% Part: lambda-calculus
% Chapter: syntax
% Section: de-bruijn

\documentclass[../../../include/open-logic-section]{subfiles}

\begin{document}

\olfileid[hi]{lam}{syn}{deb}

\olsection{डे ब्रुइन सूचकांक}

$\alpha$-तुल्यता अत्यंत स्वाभाविक है, क्योंकि $\alpha$-तुल्य पदों का
``अर्थ एक ही'' होता है। वास्तव में, लैम्ब्डा पदों का ऐसा वाक्यविन्यास देना
संभव है जो परस्पर $\alpha$-परिवर्तनीय पदों में भेद न करे। सबसे प्रसिद्ध
विधि चरों को उनके \emph{डे ब्रुइन सूचकांक} से बदल देती है।

जब हम $\lambd[x][M]$ लिखते हैं, तब स्पष्ट रूप से कहते हैं कि $x$ फलन का
प्राचल है, ताकि~$M$ में इस प्राचल के संदर्भ के लिए $x$ का उपयोग कर सकें।
किंतु डे ब्रुइन सूचकांक में प्राचलों के नाम नहीं होते; फलन के मुख्य भाग में
उनका संदर्भ ऐसी संख्या से दर्शाया जाता है जो संदर्भ और प्राचल के बीच
अमूर्तन के स्तरों की संख्या बताती है। उदाहरणार्थ,
$\lambd[x][\lambd[y][y x]]$ पर विचार करें: बाहरी अमूर्तन चर~$x$ को बाँधता
है; भीतरी अमूर्तन चर~$y$ को बाँधता है; उपपद $y x$ भीतरी अमूर्तन की
व्याप्ति में है। $y$ और उसके अमूर्त~$\lambd[y]$ के बीच कोई अमूर्तन नहीं,
पर $x$ और उसके अमूर्त~$\lambd[x]$ के बीच एक अमूर्तन है। अतः $y x$ के लिए
$0\, 1$ और पूरे पद के लिए $\lambd[][\lambd[][01]]$ लिखते हैं।

\begin{defn}
  डे ब्रुइन पद आगमनतः इस प्रकार परिभाषित हैं:
  \begin{enumerate}
  \item $n$, जहाँ $n$ कोई भी प्राकृतिक संख्या है।
  \item $PQ$, जहाँ $P$ और $Q$ दोनों डे ब्रुइन पद हैं।
  \item $\lambd[][N]$, जहाँ $N$ कोई डे ब्रुइन पद है।
  \end{enumerate}
\end{defn}

साधारण लैम्ब्डा पदों से डे ब्रुइन-सूचकांकित पदों में रूपांतरण को औपचारिक
रूप से इस प्रकार परिभाषित किया जाता है:
\begin{defn}
  \begin{align*}
    F_\Gamma(x) &= \Gamma(x) \\
    F_\Gamma(PQ) &= F_\Gamma(P)F_\Gamma(Q) \\
    F_\Gamma(\lambd[x][N]) &= \lambd[][F_{x,\Gamma}(N)]
  \end{align*}
  जहाँ $\Gamma$ शून्य से सूचकांकित चरों की सूची है और $\Gamma(x)$,
  ~$\Gamma$ में चर $x$ का स्थान दर्शाता है। उदाहरणार्थ, यदि $\Gamma$,
  ~$x,y,z$ है, तो $\Gamma(x)$, ~$0$ और $\Gamma(z)$, ~$2$ है।
  
  $x,\Gamma$ उस सूची को दर्शाता है जो $x$ को $\Gamma$ के आरंभ में जोड़ने
  से मिलती है; मसलन, पिछले उदाहरण वाला $\Gamma$ लेते हुए $w,\Gamma$,
  ~$w,x,y,z$ है।
\end{defn}

किसी डे ब्रुइन पद से मानक लैम्ब्डा पद इस प्रकार पुनः प्राप्त किया जाता है:

\begin{defn}
  \begin{align*}
    G_\Gamma(n) &= \Gamma[n] \\
    G_\Gamma(PQ) &= G_\Gamma(P) G_\Gamma(Q) \\
    G_\Gamma(\lambd[][N]) &= \lambd[x][G_{x,\Gamma}(N)]
  \end{align*}
  जहाँ $\Gamma$ फिर शून्य से सूचकांकित चरों की सूची है और $\Gamma[n]$
  स्थान~$n$ पर स्थित चर को दर्शाता है। उदाहरणार्थ, यदि $\Gamma$,
  ~$x,y,z$ है, तो $\Gamma[1]$, ~$y$ है।

  अंतिम समीकरण में चर $x$ ऐसा कोई भी चर चुना जाता है जो $\Gamma$ में न हो।
\end{defn}

यहाँ हम कुछ परिणाम बिना प्रमाण के देते हैं:

\begin{prop}
  यदि $M \aconvone M'$ और $\Gamma$ कोई ऐसी सूची है जिसमें $\FV{M}$
  सम्मिलित है, तो $F_\Gamma(M) \eqs F_\Gamma(M')$।
\end{prop}

\end{document}
